% SoftwareX (Elsevier) Original Software Publication scaffold.
% Draft planning settings; the publisher guide could not be retrieved during
% review on 2026-10-05. Confirm the current template and limits before submission:
% https://www.sciencedirect.com/journal/softwarex/publish/guide-for-authors
%   - Submissions are ONLY considered on the journal's own template; treat this
%     scaffold as a content skeleton to transfer into the official LaTeX template.
%   - Maximum 3,000 words including abstract, running text, captions and
%     footnotes (excluding title, authors, affiliations, references, metadata).
%   - Maximum 6 figures. Single column. Line numbers required.
%   - At least one of the two mandatory metadata tables must be completed.
%   - Supply the code, license, documentation, and a reproducible release link;
%     confirm the publisher's current repository requirements.
\documentclass[preprint,review,12pt]{elsarticle}

\usepackage{graphicx}
\usepackage{url}
\usepackage{amsmath}
\usepackage{amssymb}
\usepackage{lineno}

\journal{SoftwareX}

\begin{document}

\begin{frontmatter}

\title{{{TITLE}}}
\author{{{AUTHORS}}}

\begin{abstract}
{ABSTRACT} % <= 250 words, factual, self-contained, no references
\end{abstract}

\begin{keyword}
TBD. % 1-7 keywords, single words preferred
\end{keyword}

\end{frontmatter}
\linenumbers

\section{Problems and background}
{INTRO}
% The scientific problem the software addresses and why existing options are
% insufficient; the manuscript is an accompanying note to the software itself.

\section{Software description}
{METHODS}
% Architecture, main functionalities, dependencies, installation.

\subsection{Software architecture}
TBD.

\subsection{Software functionalities}
TBD.

\section{Illustrative examples}
{RESULTS}
% At least one concrete end-to-end example with inputs and expected outputs.

{FIGURES}

\section{Impact}
{DISCUSSION}
% Who benefits, what new research questions become feasible, reuse potential
% beyond the original domain; SoftwareX weighs cross-domain reuse explicitly.

\section{Conclusions}
TBD.

\section*{Declaration of competing interest}
TBD.

\section*{Acknowledgments}
TBD.

% ---------------- Mandatory metadata table ----------------
% Fill and keep in the manuscript; the official template provides the exact
% layout. At least one metadata table is mandatory.
\begin{table}[h]
\footnotesize
\begin{tabular}{|l|p{5cm}|p{6cm}|}
\hline
\textbf{Nr.} & \textbf{Code metadata description} & \textbf{Fill in} \\
\hline
C1 & Current code version & TBD. \\
C2 & Permanent link to code/repository & TBD. \\ % GitHub URL
C3 & Permanent link to reproducible capsule & TBD. \\ % e.g. Zenodo DOI
C4 & Legal code license & TBD. \\ % OSI-recognized license
C5 & Code versioning system used & git \\
C6 & Software code languages, tools, and services used & TBD. \\
C7 & Compilation requirements, operating environments and dependencies & TBD. \\
C8 & Link to developer documentation/manual & TBD. \\
C9 & Support email for questions & TBD. \\
\hline
\end{tabular}
\caption{Software metadata (mandatory).}
\end{table}

\bibliographystyle{elsarticle-num}
\bibliography{references}

\end{document}
